\chapter{Boolean Analysis — Fourier Concentration}
%%% generated by the blueprint pipeline from TCSlib.BooleanAnalysis.LMN.FourierConcentration
\section{Overview}

This module proves O'Donnell's Lemma 4.21: if a Boolean function $f$ on $n$ variables
becomes a shallow decision tree with high probability under a Bernoulli($p$)-random
restriction, then the Fourier spectrum of its $\pm1$-encoding is concentrated on low
degrees. Concretely, if $\varepsilon$ is the probability that the restricted function
requires decision-tree depth at least $k$, then the Fourier weight of $\sigma\circ f$
above degree $3k/p$ is at most $4\varepsilon$; the supporting lemmas partition tail
events over intersection patterns and aggregate the expected restricted spectrum.

%%%%%%%%%%%%%%%%%%%%%%%%%%%%%%%%%%%%%%%%%%%%%%%%%%%%%%%%%%%%%%%%%%%%%%%%%%%%%%%
\section{Declarations}

\begin{lemma}[Sign encoding of a Boolean function is $\pm1$-valued]
\lean{RestrictionFourier.isPmOne_boolToSign}
\label{RestrictionFourier.isPmOne_boolToSign}
\leanok
\uses{BooleanAnalysis.boolToSign, BooleanAnalysis.isPmOne}
\difficulty{2}
Let $g : \{0,1\}^n \to \{0,1\}$ be a Boolean-valued function, and let $\sigma$ denote
the sign encoding of a bit, $\sigma(\mathtt{false}) = 1$ and $\sigma(\mathtt{true}) =
-1$. Then the real-valued function $x \mapsto \sigma(g(x))$ takes only the values $1$
and $-1$ at every point $x$ of the cube.
\end{lemma}

\begin{lemma}[Nonnegativity of Bernoulli-restriction probabilities]
\lean{RestrictionFourier.bernoulliRestrProb_nonneg}
\label{RestrictionFourier.bernoulliRestrProb_nonneg}
\leanok
\uses{SwitchingLemma2.Restriction, SwitchingLemma2.bernoulliRestrProb, SwitchingLemma2.bernoulliRestrWeight_nonneg'}
\difficulty{2}
Let $p \in \bbr$ with $0 \le p \le 1$, and let $E$ be a decidable event on restrictions
of $n$ variables, where a restriction assigns each coordinate either a fixed Boolean
value or leaves it free. Under the Bernoulli($p$)-random restriction --- each
restriction $\rho$ with free-variable set $J$ carries weight
$p^{\abs{J}}\bigl(\tfrac{1-p}{2}\bigr)^{\,n-\abs{J}}$ --- the probability of $E$,
namely the weighted sum of the indicator of $E$ over all restrictions, is nonnegative:
$0 \le \Pr_\rho[E]$.
\end{lemma}

\begin{lemma}[Partitioning a tail event over intersection patterns]
\lean{RestrictionFourier.sum_prob_inter_eq_card_ge}
\label{RestrictionFourier.sum_prob_inter_eq_card_ge}
\leanok
\uses{SwitchingLemma2.Restriction, SwitchingLemma2.Restriction.freeVars, SwitchingLemma2.bernoulliRestrProb, SwitchingLemma2.bernoulliRestrWeight}
\difficulty{4}
Fix $n$, a real parameter $p$, a subset $U \subseteq [n]$, and a threshold $k \in \bbn$.
For a restriction $\rho$ on $n$ variables write $J$ for its set of free variables, and
write $\Pr_\rho[\cdot]$ for the sum, over all restrictions, of the Bernoulli weight
$p^{\abs{J}}\bigl(\tfrac{1-p}{2}\bigr)^{\,n-\abs{J}}$ times the indicator of the event
(a genuine probability when $0 \le p \le 1$). Then
\[
  \sum_{\substack{S \subseteq [n] \\ \abs{S} \ge k}}
    \Pr_\rho\bigl[\,U \cap J = S\,\bigr]
  \;=\;
  \Pr_\rho\bigl[\,\abs{U \cap J} \ge k\,\bigr].
\]
The identity holds for every real $p$; no bound on $p$ is assumed.
\end{lemma}

\begin{lemma}[Tail of the expected restricted Fourier spectrum]
\lean{RestrictionFourier.sum_tail_expectation_eq}
\label{RestrictionFourier.sum_tail_expectation_eq}
\leanok
\uses{BooleanAnalysis.BooleanFunc, BooleanAnalysis.fourierCoeff, RestrictionFourier.bernoulliRestrProb_inter_freeVars, RestrictionFourier.expectation_fourierCoeff_sq_restrictBF, RestrictionFourier.restrictBF, RestrictionFourier.sum_prob_inter_eq_card_ge, SwitchingLemma2.Restriction, SwitchingLemma2.Restriction.freeVars, SwitchingLemma2.bernoulliRestrProb, SwitchingLemma2.bernoulliRestrWeight}
\difficulty{4}
Fix $n$, a real parameter $p$, a function $F : \{0,1\}^n \to \bbr$, and $k \in \bbn$.
For a restriction $\rho$ on $n$ variables with free-variable set $J$, write $w_p(\rho)
= p^{\abs{J}}\bigl(\tfrac{1-p}{2}\bigr)^{\,n-\abs{J}}$ for its Bernoulli weight,
$F_\rho$ for the restriction of $F$ by $\rho$ (free coordinates read from the input,
fixed coordinates from $\rho$), $\hat g(T)$ for the Fourier--Walsh coefficient of a
function $g$ at a set $T$, and $\Pr_\rho[\cdot]$ for the $w_p$-weighted sum of an
event's indicator over all restrictions. Then
\[
  \sum_{\substack{S \subseteq [n] \\ \abs{S} \ge k}}
    \sum_{\rho} w_p(\rho)\, \widehat{F_\rho}(S)^2
  \;=\;
  \sum_{U \subseteq [n]} \hat F(U)^2 \;
    \Pr_\rho\bigl[\,\abs{U \cap J} \ge k\,\bigr].
\]
When $0 \le p \le 1$ the left-hand side is
$\sum_{\abs{S} \ge k} \E_\rho\bigl[\widehat{F_\rho}(S)^2\bigr]$, so this is O'Donnell's
Proposition 4.17 summed over the high-degree frequency sets.
\end{lemma}

\begin{theorem}[Fourier concentration from random restrictions]
\lean{RestrictionFourier.odonnell_lemma_4_21}
\label{RestrictionFourier.odonnell_lemma_4_21}
\leanok
\uses{BooleanAnalysis.boolToSign, BooleanAnalysis.fourierCoeff, BooleanAnalysis.parseval_pm_one, DecisionTree.degree_le_dtDepth, RestrictionFourier.bernoulliRestrProb_card_inter_ge, RestrictionFourier.bernoulliRestrProb_nonneg, RestrictionFourier.isPmOne_boolToSign, RestrictionFourier.restrictBF, RestrictionFourier.restrictBF_boolToSign, RestrictionFourier.sum_tail_expectation_eq, SwitchingLemma2.Restriction, SwitchingLemma2.Restriction.freeVars, SwitchingLemma2.bernoulliRestrProb, SwitchingLemma2.bernoulliRestrWeight, SwitchingLemma2.bernoulliRestrWeight_nonneg', SwitchingLemma2.restrictFn, dtDepth}
\difficulty{5}
This is O'Donnell's Lemma 4.21, with constant $4$ in place of $3$. Let $f : \{0,1\}^n
\to \{0,1\}$ be a Boolean function, let $p \in \bbr$ with $0 \le p \le 1$, and let $k
\in \bbn$ with $k \ge 1$. Let $\rho$ be a Bernoulli($p$)-random restriction on $n$
variables: each coordinate is independently left free with probability $p$ and
otherwise fixed to a uniformly random bit, so a restriction with free-variable set $J$
has probability $p^{\abs{J}}\bigl(\tfrac{1-p}{2}\bigr)^{\,n-\abs{J}}$. Write $f_\rho$
for the function obtained from $f$ by fixing the coordinates outside $J$ according to
$\rho$, $\mathrm{DT}(g)$ for the least depth of a decision tree computing a Boolean
function $g$, and $\sigma$ for the sign encoding ($\sigma(\mathtt{false}) = 1$,
$\sigma(\mathtt{true}) = -1$). Setting
$\varepsilon = \Pr_\rho\bigl[\mathrm{DT}(f_\rho) \ge k\bigr]$, the Fourier spectrum of
the $\pm1$-encoding $\sigma \circ f$ is $4\varepsilon$-concentrated below degree
$3k/p$:
\[
  \sum_{\substack{U \subseteq [n] \\ 3k \le p\,\abs{U}}}
    \widehat{\sigma \circ f}(U)^2 \;\le\; 4\,\varepsilon,
\]
where the degree condition is stated division-free as $3k \le p\,\abs{U}$.
\end{theorem}
